\documentclass[letterpaper]{article}%
\usepackage{lastpage}%
\usepackage{parskip}%
\usepackage{geometry}%
\geometry{left=22.5mm,right=16.1mm,top=48mm,bottom=25mm,headheight=12.5mm,footskip=12.5mm}%
\usepackage{fontspec}%
\usepackage[spanish,activeacute]{babel}%
\usepackage{anyfontsize}%
\usepackage{fancyhdr}%
\usepackage{csquotes}%
\usepackage[style=ieee,backend=biber]{biblatex}%
\usepackage[skins,breakable]{tcolorbox}%
\usepackage{color}%
\usepackage{tikz}%
\usepackage{tabularx}%
\usepackage{longtable}%
\usepackage{ltablex}%
\usepackage{ragged2e}%
\usepackage{graphicx}%
\usepackage{xcolor}%
%
\setmainfont{Arial}%
\addbibresource{../bibliografia.bib}%
\renewcommand*{\bibfont}{\fontsize{12}{16}\selectfont}%
\definecolor{gris}{rgb}{0.27,0.27,0.27}%
\definecolor{parte}{rgb}{0.02,0.204,0.404}%
\definecolor{azulsuaveTEC}{rgb}{0.02,0.455,0.773}%
\definecolor{fila}{rgb}{0.929,0.929,0.929}%
\definecolor{linea}{rgb}{0.749,0.749,0.749}%
\fancypagestyle{headfoot}{%
\renewcommand{\headrulewidth}{0pt}%
\renewcommand{\footrulewidth}{0pt}%
\fancyhead{%
}%
\fancyfoot{%
}%
\fancyhead[L]{%
\begin{minipage}{0.5\textwidth}%
\flushleft%
\includegraphics[width=62.5mm]{../figuras/Logo.png}%
\end{minipage}%
}%
\fancyfoot[L]{%
\textcolor{azulsuaveTEC}{%
Escuela de Ingeniería Electromecánica%
}%
\par \parbox{0.85\textwidth}{%
\fontsize{8}{0}\selectfont \textcolor{azulsuaveTEC}{Carreras de: Ingeniería en Mantenimiento Industrial; Ingeniería en Electrónica; Ingeniería en Producción Industrial; Ingeniería Mecatrónica e Ingeniería en Materiales}%
}%
}%
\fancyfoot[R]{%
\textcolor{azulsuaveTEC}{%
Página \thepage \hspace{1pt} de \pageref*{LastPage}%
}%
}%
}%
%
\begin{document}%
\normalsize%
\thispagestyle{empty}%
\begin{tikzpicture}[overlay,remember picture]%
\node[inner sep = 0mm,outer sep = 0mm,anchor = north west,xshift = -23mm,yshift = 22mm] at (0.0,0.0) {\includegraphics[width=21cm]{../figuras/Logo_portada.png}};%
\end{tikzpicture}%
\vspace*{150mm}%
\par\fontsize{14}{0}\selectfont \textcolor{black}{Programa del curso IEM2101}%
\par\fontsize{18}{25}\selectfont \textbf{\textcolor{azulsuaveTEC}{Dibujo técnico asistido por computadora}}%
\begin{tabularx}{\textwidth}{m{0.02\textwidth}m{0.98\textwidth}}%
&\hspace*{0mm}\fontsize{12}{14}\selectfont \textbf{\textcolor{gris}{Escuela de Ingeniería Electromecánica}}\\%
[-4pt]%
&\hspace*{0mm}\fontsize{12}{14}\selectfont \textbf{\textcolor{gris}{Carreras de: Ingeniería en Mantenimiento Industrial; Ingeniería en Electrónica; Ingeniería en Producción Industrial; Ingeniería Mecatrónica e Ingeniería en Materiales}}\\%
\end{tabularx}%
\newpage%
\pagestyle{headfoot}%
\par\fontsize{14}{0}\selectfont \textbf{\textcolor{parte}{I parte: Aspectos relativos al plan de estudios}}%
\par\hspace*{4mm}\fontsize{12}{20}\selectfont \textbf{\textcolor{parte}{1 Datos generales}}%
\begin{tabularx}{\textwidth}{p{7cm}p{9cm}}%
\textbf{Nombre del curso:}&Dibujo técnico asistido por computadora\\%
[10pt]%
\textbf{Código:}&IEM2101\\%
[10pt]%
\textbf{Tipo de curso:}&Teórico\\%
[10pt]%
\textbf{Electivo o no:}&No\\%
[10pt]%
\textbf{Nº de créditos:}&3\\%
[10pt]%
\textbf{Nº horas de clase por semana:}&4\\%
[10pt]%
\textbf{Nº horas extraclase por semana:}&5\\%
[10pt]%
\textbf{\% de areas curriculares:}&6.0\% del area: \textbf{mecánica}\\%
[10pt]%
\textbf{Ubicación en el plan de estudios:}&Curso de 1
 semestre en Ingeniería en Electrónica e Ingeniería en Producción Industrial. Curso de 2
 semestre en Ingeniería en Mantenimiento Industrial. Curso de 3
 semestre en Ingeniería Mecatrónica. Curso de 4\textsuperscript{to} semestre en Ingeniería en Materiales. \\%
[10pt]%
\textbf{Requisitos:}&nan\\%
[10pt]%
\textbf{Correquisitos:}&nan\\%
[10pt]%
\textbf{El curso es requisito de:}&nan\\%
[10pt]%
\textbf{Asistencia:}&nan\\%
[10pt]%
\textbf{Suficiencia:}&nan\\%
[10pt]%
\textbf{Posibilidad de reconocimiento:}&nan\\%
[10pt]%
\textbf{Vigencia del programa:}&nan\\%
[10pt]%
\end{tabularx}%

\begin{tcolorbox}[
blanker,
width=0.78\textwidth,enlarge left by=0.24\textwidth,
before skip=6pt,
breakable,
overlay unbroken and first={%
    \node[inner sep=0pt,outer sep=0pt,text width=0.22\textwidth,
    align=none,
    below right]
    at ([xshift=-0.24\textwidth]frame.north west)
{
\hspace*{0mm}\fontsize{12}{14}\selectfont \textbf{\textcolor{parte}{2 Descripción general}}
};}]
El curso de Dibujo Técnico Asistido por Computadora contribuye significativamente a la formación y desarrollo profesional de los estudiantes, equipándolos con las habilidades y herramientas necesarias para la comunicación, el diseño y la ejecución en el ámbito de la ingeniería.

Entre los aprendizajes más destacados se encuentran:  interpretar y aplicar las normas INTE/ISO de dibujo técnico en situaciones prácticas; visualizar y representar objetos tridimensionales en un plano bidimensional y viceversa; desarrollar destrezas en el uso de herramientas CAD para la elaboración de planos; y además de las habilidades técnicas, se busca promover el compromiso, el respeto y la ética profesional entre los participantes.​

Este curso: sienta las bases fundamentales para asignaturas más avanzadas en el campo del diseño y la ingeniería mecánica. Proporciona una comprensión sólida de los principios y técnicas de dibujo técnico que son esenciales para el desarrollo de proyectos más complejos las áreas de diseño y manufactura. De esta manera, establece una conexión con otros cursos de la carrera, preparando a los estudiantes para enfrentar desafíos en su trayectoria académica y profesional.​

En caso de que un estudiante requiera apoyos educativos, podrá solicitarlos a través del Departamento de Orientación y Psicología. 
\end{tcolorbox}
%
\vspace*{10mm}%

\begin{tcolorbox}[
blanker,
width=0.78\textwidth,enlarge left by=0.24\textwidth,
before skip=6pt,
breakable,
overlay unbroken and first={%
    \node[inner sep=0pt,outer sep=0pt,text width=0.22\textwidth,
    align=none,
    below right]
    at ([xshift=-0.24\textwidth]frame.north west)
{
\hspace*{0mm}\fontsize{12}{14}\selectfont \textbf{\textcolor{parte}{3 Objetivos}}
};}]
Al final del curso la persona estudiante será capaz de realizar un plano de una pieza mecánica, según las normas INTE/ISO, que contenga información necesaria y suficiente para la interpretación de la forma y dimensiones de la pieza.\\
La persona estudiante será capaz de:\\
\hspace*{0.02\linewidth}\parbox{0.98\linewidth}{\strut\textbullet\, Aplicar las normas INTE/ISO y de otros estándares.\strut}\\
\hspace*{0.02\linewidth}\parbox{0.98\linewidth}{\strut\textbullet\, Representar mediante proyecciones ortogonales, proyecciones axonométricas, vistas auxiliares y vistas de corte una pieza mecánica.\strut}\\
\hspace*{0.02\linewidth}\parbox{0.98\linewidth}{\strut\textbullet\, Acotar adecuadamente una pieza mecánica en las proyecciones ortogonales.\strut}\\
 
\end{tcolorbox}
%
\vspace*{10mm}%

\begin{tcolorbox}[
blanker,
width=0.78\textwidth,enlarge left by=0.24\textwidth,
before skip=6pt,
breakable,
overlay unbroken and first={%
    \node[inner sep=0pt,outer sep=0pt,text width=0.22\textwidth,
    align=none,
    below right]
    at ([xshift=-0.24\textwidth]frame.north west)
{
\hspace*{0mm}\fontsize{12}{14}\selectfont \textbf{\textcolor{parte}{4 Contenidos}}
};}]
1. Generalidades\\
\hspace*{0.02\linewidth}\parbox{0.98\linewidth}{\strut 1.1. El Dibujo Técnico como lenguaje. Historia del Dibujo Técnico.\strut}\\
\hspace*{0.02\linewidth}\parbox{0.98\linewidth}{\strut 1.2. Objetivos del curso\strut}\\
\hspace*{0.02\linewidth}\parbox{0.98\linewidth}{\strut 1.3. Instrumentos de Dibujo: Escalímetro, escuadras.\strut}\\
\hspace*{0.02\linewidth}\parbox{0.98\linewidth}{\strut 1.4. Rotulado técnico y Formatos para dibujo Técnico\strut}\\
\hspace*{0.02\linewidth}\parbox{0.98\linewidth}{\strut 1.5. Norma de rotulado INTE-ISO 3098/0/2/3-2008\strut}\\
\hspace*{0.02\linewidth}\parbox{0.98\linewidth}{\strut 1.6. Formatos según INTE-ISO 5457-2008\strut}\\
\hspace*{0.02\linewidth}\parbox{0.98\linewidth}{\strut 1.7. Información que debe contener un cajetín INTE-ISO 7200-2008\strut}\\
2. Geometría Descriptiva (2 hrs)\\
\hspace*{0.02\linewidth}\parbox{0.98\linewidth}{\strut 2.1. Consideraciones fundamentales de la Geometría Descriptiva :\strut}\\
\hspace*{0.04\linewidth}\parbox{0.96\linewidth}{\strut 2.1.1. Objetivos del curso\strut}\\
\hspace*{0.04\linewidth}\parbox{0.96\linewidth}{\strut 2.1.2. Concepto de Geometría Descriptiva y su Historia\strut}\\
\hspace*{0.02\linewidth}\parbox{0.98\linewidth}{\strut 2.2. Proyección del punto, el segmento y los planos en el espacio\strut}\\
\hspace*{0.04\linewidth}\parbox{0.96\linewidth}{\strut 2.2.1. Proyección de un punto, el segmento y los planos en las vistas\strut}\\
\hspace*{0.04\linewidth}\parbox{0.96\linewidth}{\strut 2.2.2. Proyección de un punto, el segmento y los planos en el espacio\strut}\\
\hspace*{0.02\linewidth}\parbox{0.98\linewidth}{\strut 2.3. Longitudes y dimensiones naturales\strut}\\
\hspace*{0.04\linewidth}\parbox{0.96\linewidth}{\strut 2.3.1. Procedimiento de rotación para encontrar la dimensión real de un segmento y un plano;\strut}\\
\hspace*{0.04\linewidth}\parbox{0.96\linewidth}{\strut 2.3.2. Procedimiento de sustitución de planos para encontrar la dimensión real de un segmento y un plano;\strut}\\
\hspace*{0.04\linewidth}\parbox{0.96\linewidth}{\strut 2.3.3. Procedimiento de superposición para encontrar la dimensión real de un segmento y un plano.\strut}\\
\hspace*{0.02\linewidth}\parbox{0.98\linewidth}{\strut 2.4. Características particulares de la representación de los Cuerpos Geométricos\strut}\\
\hspace*{0.04\linewidth}\parbox{0.96\linewidth}{\strut 2.4.1. Características particulares del prisma y su representación;\strut}\\
\hspace*{0.04\linewidth}\parbox{0.96\linewidth}{\strut 2.4.2. Características del cono y su representación;\strut}\\
\hspace*{0.04\linewidth}\parbox{0.96\linewidth}{\strut 2.4.3. Características de la pirámide y su representación;\strut}\\
\hspace*{0.04\linewidth}\parbox{0.96\linewidth}{\strut 2.4.4. Características del cilindro y su representación;\strut}\\
\hspace*{0.04\linewidth}\parbox{0.96\linewidth}{\strut 2.4.5. Características del toroide y su representación;\strut}\\
\hspace*{0.04\linewidth}\parbox{0.96\linewidth}{\strut 2.4.6. Características de la esfera y su representación\strut}\\
3. Escalas (4 hrs)\\
\hspace*{0.02\linewidth}\parbox{0.98\linewidth}{\strut 3.1. Concepto de escalas\strut}\\
\hspace*{0.02\linewidth}\parbox{0.98\linewidth}{\strut 3.2. Escalas según INTE ISO 5455-2008\strut}\\
4. Proyecciones ortogonales (14 hrs)\\
\hspace*{0.02\linewidth}\parbox{0.98\linewidth}{\strut 4.1. Sistema de proyección y designación de vistas según la norma INTE- ISO 128/30/34-2008\strut}\\
\hspace*{0.02\linewidth}\parbox{0.98\linewidth}{\strut 4.2. Criterios de selección de la vista frontal y la ubicación de las otras vistas.\strut}\\
\hspace*{0.02\linewidth}\parbox{0.98\linewidth}{\strut 4.3. Cantidad de vistas que definen un objeto.\strut}\\
\hspace*{0.02\linewidth}\parbox{0.98\linewidth}{\strut 4.4. Significado y utilización de los tipos de líneas.\strut}\\
\hspace*{0.02\linewidth}\parbox{0.98\linewidth}{\strut 4.5. Tipos de líneas según la norma INTE ISO 128/20/21/22/23/24-2008\strut}\\
5. Presentación de un plano:\\
\hspace*{0.02\linewidth}\parbox{0.98\linewidth}{\strut 5.1. Calidad de líneas\strut}\\
\hspace*{0.02\linewidth}\parbox{0.98\linewidth}{\strut 5.2. Orden y adecuada ubicación de la información\strut}\\
\hspace*{0.02\linewidth}\parbox{0.98\linewidth}{\strut 5.3. Especificaciones técnicas.\strut}\\
\hspace*{0.02\linewidth}\parbox{0.98\linewidth}{\strut 5.4. Práctica de proyecciones ortogonales croquizando “a mano alzada”.\strut}\\
\hspace*{0.02\linewidth}\parbox{0.98\linewidth}{\strut 5.5. Práctica de proyecciones ortogonales utilizando Software de dibujo.\strut}\\
6. Proyecciones Axonométricas (12 hrs)\\
\hspace*{0.02\linewidth}\parbox{0.98\linewidth}{\strut 6.1. Ejes de proyección\strut}\\
\hspace*{0.02\linewidth}\parbox{0.98\linewidth}{\strut 6.2. Tipos de axonometrías\strut}\\
\hspace*{0.02\linewidth}\parbox{0.98\linewidth}{\strut 6.3. Proyecciones axonométricas a mano alzada y utilizando software de dibujo.\strut}\\
7. Acotado (4 hrs) \\
\hspace*{0.02\linewidth}\parbox{0.98\linewidth}{\strut 7.1. Normas y recomendaciones INTE/ISO 129/1-2008 sobre acotado.\strut}\\
\hspace*{0.02\linewidth}\parbox{0.98\linewidth}{\strut 7.2. Líneas utilizadas en el acotado\strut}\\
\hspace*{0.02\linewidth}\parbox{0.98\linewidth}{\strut 7.3. Posición de la cota\strut}\\
\hspace*{0.02\linewidth}\parbox{0.98\linewidth}{\strut 7.4. Rotulado de cotas\strut}\\
\hspace*{0.02\linewidth}\parbox{0.98\linewidth}{\strut 7.5. Criterios para la acotación correcta de piezas.\strut}\\
\hspace*{0.02\linewidth}\parbox{0.98\linewidth}{\strut 7.6. Relación entre cota y escala.\strut}\\
8. Cortes y secciones (4 hrs)\\
\hspace*{0.02\linewidth}\parbox{0.98\linewidth}{\strut 8.1. Concepto de cortes y secciones. Conveniencia de su utilización\strut}\\
\hspace*{0.02\linewidth}\parbox{0.98\linewidth}{\strut 8.2. Representación e indicación de cortes según la norma INTE/ISO 128/40-44- 2008\strut}\\
\hspace*{0.02\linewidth}\parbox{0.98\linewidth}{\strut 8.3. Achurado\strut}\\
\hspace*{0.02\linewidth}\parbox{0.98\linewidth}{\strut 8.4. Tipos de cortes.\strut}\\
\hspace*{0.04\linewidth}\parbox{0.96\linewidth}{\strut 8.4.1. Sección en un plano de corte.\strut}\\
\hspace*{0.04\linewidth}\parbox{0.96\linewidth}{\strut 8.4.2. Sección en dos planos paralelos.\strut}\\
\hspace*{0.04\linewidth}\parbox{0.96\linewidth}{\strut 8.4.3. Sección en tres planos de corte continuos\strut}\\
\hspace*{0.04\linewidth}\parbox{0.96\linewidth}{\strut 8.4.4. Sección en dos planos de intersección\strut}\\
\hspace*{0.04\linewidth}\parbox{0.96\linewidth}{\strut 8.4.5. Plano de corte posesionado parcialmente fuera de la pieza\strut}\\
\hspace*{0.04\linewidth}\parbox{0.96\linewidth}{\strut 8.4.6. Sección removida de una vista\strut}\\
\hspace*{0.04\linewidth}\parbox{0.96\linewidth}{\strut 8.4.7. Secciones sucesivas\strut}\\
\hspace*{0.04\linewidth}\parbox{0.96\linewidth}{\strut 8.4.8. Cortes oblicuos o auxiliares\strut}\\
\hspace*{0.04\linewidth}\parbox{0.96\linewidth}{\strut 8.4.9. Cortes y secciones parciales\strut}\\
\hspace*{0.04\linewidth}\parbox{0.96\linewidth}{\strut 8.4.10. Corte y secciones de piezas simétricas\strut}\\
9. Vistas auxiliares simples (4 hrs)\\
\hspace*{0.02\linewidth}\parbox{0.98\linewidth}{\strut 9.1. Vistas auxiliares\strut}\\
\hspace*{0.02\linewidth}\parbox{0.98\linewidth}{\strut 9.2. Tipos de vistas auxiliares\strut}\\
\hspace*{0.02\linewidth}\parbox{0.98\linewidth}{\strut 9.3. Ubicación de las vistas auxiliares\strut}\\
\hspace*{0.02\linewidth}\parbox{0.98\linewidth}{\strut 9.4. Rotulado de vistas auxiliares.\strut}\\ 
\end{tcolorbox}
%
\vspace*{10mm}%
\newpage%
\par\fontsize{14}{0}\selectfont \textbf{\textcolor{parte}{II parte: Aspectos operativos}}%
\par\fontsize{12}{0}\selectfont \textbf{\textcolor{parte}{ }}%
\renewcommand{\arraystretch}{1.5}%
\begin{tabularx}{\textwidth}{>{\raggedright}p{0.18\textwidth}p{0.72\textwidth}}%
\par\fontsize{12}{0}\selectfont \textbf{\textcolor{parte}{5 Metodología de enseñanza y aprendizaje}}&Modalidad presencial en el horario de clase correspondiente. En las sesiones se expondrán los conceptos teóricos relevantes de cada tema por medio de material escrito o audiovisual, los cuales serán trabajados por el estudiante en numerosas prácticas realizadas en clases y tareas coordinadas.\\%
\end{tabularx}%
\begin{tabularx}{\textwidth}{>{\raggedright}m{0.18\textwidth}m{0.07\textwidth}m{0.17\textwidth}m{0.17\textwidth}m{0.17\textwidth}m{0.04\textwidth}}%
\par\fontsize{12}{0}\selectfont \textbf{\textcolor{parte}{6 Evaluación}}&&\par\fontsize{12}{16}\selectfont \textbf{\textcolor{black}{Tipo}}&\par\fontsize{12}{16}\selectfont \textbf{\textcolor{black}{Cantidad}}&\par\fontsize{12}{16}\selectfont \textbf{\textcolor{black}{Porcentaje}}&\\%
[12pt]%
&&Examen Parcial&1&25
&\\%
[12pt]%
&&Examen Parcial&1&35
&\\%
[12pt]%
&&Proyecto grupal&1&25
&\\%
[12pt]%
&&Tareas digitales coordinadas&1&15&\\%
[12pt]%
\end{tabularx}%

\begin{tcolorbox}[
blanker,
width=0.78\textwidth,enlarge left by=0.24\textwidth,
before skip=6pt,
breakable,
overlay unbroken and first={%
    \node[inner sep=0pt,outer sep=0pt,text width=0.22\textwidth,
    align=none,
    below right]
    at ([xshift=-0.24\textwidth]frame.north west)
{
\hspace*{0mm}\fontsize{12}{14}\selectfont \textbf{\textcolor{parte}{7 Bibliografía}}
};}]

\nocite{hernandez2013dibujo}
\nocite{lieu2011dibujo}
\nocite{bogoliubov1988dibujo}
\nocite{bogoliubov1989tareas}
\printbibliography[heading=none] 
\end{tcolorbox}
%
\vspace*{10mm}%

\begin{tcolorbox}[
blanker,
width=0.78\textwidth,enlarge left by=0.24\textwidth,
before skip=6pt,
breakable,
overlay unbroken and first={%
    \node[inner sep=0pt,outer sep=0pt,text width=0.22\textwidth,
    align=none,
    below right]
    at ([xshift=-0.24\textwidth]frame.north west)
{
\hspace*{0mm}\fontsize{12}{14}\selectfont \textbf{\textcolor{parte}{8 Profesor}}
};}]
M.Sc. Luis Felipe Córdoba Ramírez\\
Máster en ciencias en Ingeniería Mecánica. Universidad de Utah. Estados Unidos.\\
Licenciado en Ingeniería en Mantenimiento Industrial. Instituto Tecnológico de Costa Rica. Costa Rica\\
Correo: lfcordoba@itcr.ac.cr. Oficina: 22\\
Escuela de Ingeniería Electromecánica. Cartago\\[12pt]
Christopher Vega Sánchez Ph.D.\\
Doctor en filosofía en ciencias. Universidad de Sídney, Australia. \\
Máster en ciencias en Ingeniería de Sistemas Microelectromecánicos. Universidad de Freiburg. Alemania.\\
Licenciado en Ingeniería en Mantenimiento Industrial,Instituto Tecnológico de Costa Rica. Costa Rica.\\
Correo: cvega@itcr.ac.cr. Oficina: 20\\
Escuela de Ingeniería Electromecánica. Cartago\\[12pt]
 
\end{tcolorbox}
%
\end{document}